\begin{abstract}

Existing Text-to-Image (T2I) safety benchmarks rely on binary safe/unsafe classification, which fails to capture the ``safety gap'' in which subtly escalated prompts pass text filters yet produce harmful visuals.
We present \textbf{\dataname}, a benchmark for boundary-aware safety evaluation built from 1,026 adversarial severity ladders spanning 11 risk categories.
Each ladder escalates through four rungs (Safe, Low-risk, Mid-risk, High-risk) with paired prompts, reference images, and risk explanations, all constructed automatically by our Adversarial Severity Ladders (ASL) pipeline without manual authoring.
Using \dataname, we evaluate 6 T2I models, 9 defense methods, 5 safety filters, and 5 adversarial attack strategies, yielding 7 insights into the T2I safety ecosystem.
Our central finding is a universal safety gap: existing mechanisms are calibrated for high-risk (L3) violations while leaving the low-to-mid-risk (L1--L2) regime largely undefended, a failure mode invisible to binary benchmarks.
A severity-aware filter fine-tuned on ladder data attains \red{72.8\%} detection at L1 with a \red{0.5\%} L0 false-positive rate, cutting the overall bypass rate from 88.0\% to \red{25.2\%} and \red{dominating all five baselines on every metric evaluated}, showing that \dataname closes the gap it diagnoses.
\dataname and the ASL pipeline will be publicly released to support boundary-aware safety research.

\end{abstract}
